\section{Findings}
\subsection{Welch t-test: The Entry Shock Analysis (Section 1.3)}
The Welch t-tests quantify the "Architectural Footprint" of a developer's first interactions.
\begin{itemize}
\item \textbf{Initial Breadth:} At the first commit, Founders touch significantly more packages ($p=0.0067$) and classes ($p=0.0042$) than Late Joiners. This indicates a ``Blank Canvas'' phase where founders must define structure, whereas Late Joiners are immediately constrained by established abstractions.
\item \textbf{Discovery Convergence:} By 10 commits, the difference in package touches between Founders and Late Joiners becomes non-significant ($p=0.3462$ for Moderate, $p=0.4084$ for Sustained), suggesting that high-level architectural discovery stabilizes quickly regardless of entry timing.
\end{itemize}  
\subsection{Two-Way ANOVA: Role vs. Tenure (Section 1.4)}
The ANOVA results provide deep insight into the drivers of developer behavior.
\begin{itemize}
\item \textbf{Main Effects:} The significant Role effect for classes ($p=0.0000$) and methods ($p=0.0088$) at the first commit suggests that \textit{when} a developer joins is the primary determinant of early exploration depth.
\item \textbf{Interaction Effects:} Significant interaction effects ($p=0.0000$ for classes, $p=0.0081$ for methods at commit 1) demonstrate that the behavior of a Founder who eventually becomes Sustained is fundamentally distinct from a Late Joiner who achieves the same status. This confirms that entry timing and future engagement are coupled variables in codebase learning.
\end{itemize}  
\subsection{Tukey HSD: Behavioral Divergence (Section 1.5)}
Post-hoc analysis identifies the specific points where the learning curve deviates.
\begin{itemize}
\item \textbf{The Retention Signal:} A critical divergence is found between \textbf{Late Joiner Moderate} and \textbf{Late Joiner Sustained} groups at the first commit for classes ($p=0.0000$) and methods ($p=0.0000$).
\item \textbf{Analysis:} This suggests that the very first contribution acts as a proxy for future engagement; those destined to stay (Sustained) interact with the codebase with significantly higher intensity from day one compared to those who will eventually churn (Moderate).
\end{itemize}  \subsection{Logistic Regression: Predictive Modeling (Section 1.6)}
The regression model converts behavioral observations into a predictive tool for project maintainers.
\begin{itemize}
\item \textbf{Effect Size:} The positive coefficients and Odds Ratios $>1$ across all components (packages, classes, methods) confirm that increased architectural engagement in early commits is a statistically significant predictor of long-term retention.
\item \textbf{Identification Thresholds:} The model defines a 50\% Probability Threshold'' for becoming a Sustained developer  . For instance, a developer touching $\sim 18$ packages within 10 commits has a 50\% chance of becoming a Sustained contributor  . This identifies a Sweet Spot'' of exploration that correlates with high project stickiness.
\end{itemize}  
\section{Overall Summary of Findings}
The research establishes that codebase learning is not a uniform process but is heavily influenced by entry timing and early interaction intensity. Founders act as architects, touching multiple components to build the system, while Late Joiners must overcome a higher cognitive barrier to achieve the same breadth of knowledge. Crucially, the intensity of the very first commit—measured by unique class and method touches—serves as a statistically significant indicator of whether a developer will transition from a transient contributor to a sustained project member[cite: 5]. This suggests that the complexity and scope of initial tasks are vital levers for improving developer retention in open-source ecosystems[cite: 5].  